Text for Kets

 One consequence of this deserves particular notice, since it will concern us a good deal in our classification of the sciences, and yet is quite usually overlooked and assumed not to be as it is. Namely, it follows that it may be quite impossible to draw a sharp line of demarcation between two classes, although they are real and natural classes in strictest truth. Namely, this will happen when the form about which the individuals of one class cluster is not so unlike the form about which individuals of another class cluster but that variations from each middling form may precisely agree. In such a case, we may know in regard to any intermediate form what proportion of the objects of that form had one purpose and what proportion the other; but unless we have some supplementary information we cannot tell which ones had one purpose and which the other.
209. The reader may be disposed to suspect that this is merely a mathematician's fancy, and that no such case would be likely ever to occur. But he may be assured that such occurrences are far from being rare. In order to satisfy him that this state of things does occur, I will mention an incontestable instance of it; — incontestable, at least, by any fair mind competent to deal with the problem. Prof. [W. M.] Flinders Petrie, whose reasoning powers I had admired long before his other great scientific qualities had been proved, among which his great exactitude and circumspection as a metrologist concerns us here, exhumed, at the ancient trading town of Naucratis, no less than 158 balance-weights having the Egyptian ket as their unit.•P1 The great majority of them are of basalt and syenite, material so unchangeable that the corrections needed to bring them to their original values are small. I shall deal only with 144 of them from each of which Mr. Petrie has calculated the value of the ket to a tenth of a Troy grain. Since these values range all the way from 137 to 152 grains, it is evident that the weights were intended to be copies of several different standards, probably four or five; for there would be no use of a balance, if one could detect the errors of the balance-weights by simply »hefting« them, and comparing them with one's memory of the standard weight. Considering that these weights are small, and were therefore used for weighing costly or even precious matter, our knowledge of the practice of weighing among the ancients gives us ground for thinking it likely that about half the weights would depart from their virtual standards by more, and about half by less, than, say, four or five tenths of one per cent, which, upon a ket, would be from half to two-thirds of a grain. Now the whole interval here is fourteen and one-half grains; and between 136.8 grains to 151.3 grains there is no case of an interval of more than a third of a grain not represented by any weight among the 144. To a person thoroughly familiar with the theory of errors this shows that there must be four or five different standards to which different ones aim to conform. . . . In order to represent these observations, I have adopted the following rough-and-ready theory; for to make elaborate calculations would, from every point of view, be a waste of time. I have assumed that there were five different standards; that the weights depart from their standards according to the probability curve; and that the probable error of a single weight is five-eighths of a grain. I assume that of the 144 weights
36 were designed to conform to a standard of 139.2 grs.
25 were designed to conform to a standard of 142.2 grs.
26 were designed to conform to a standard of 144.7 grs.
23 were designed to conform to a standard of 146.95 grs.
34 were designed to conform to a standard of 149.7 grs.
. . . I repeat that this theory has not been the subject of any but the simplest calculations. It is obvious that some such theory must be true; but to decide how near my theory probably comes to the true theory or how it ought to be modified, would be a very intricate problem for the solution of which the data are probably insufficient. It does not concern us here; our object being merely to make it clear that truly natural classes may, and undoubtedly often do merge into one another inextricably.
It is, I think, pretty certain that there were as many as five standards. Before the adoption of the metric system, every city throughout the greater part, if not all, the continent of Europe had its own pound, like its own patois. See the article »pound« in the Century Dictionary,1) which was based on a list of some three hundred of such pounds whose values were known to me, a list now kept in manuscript in the Astor Library. That the same state of things must have been true in ancient Egypt may be inferred from the looseness of the tie which bound the different provinces of that empire together. Even their religions were different; so that a fortiori their kets would be so. Besides, none of the kets carry any authoritative mark; which is pretty conclusive proof that the central government did not intervene. It is, therefore, probable that the five standards were those of five towns with which Naucratis carried on trade. Yet virtual standards may be created in other ways. For example, where government does not insure uniformity in weights, it is usual for buyers to bring their own weights. It would thus naturally happen that some balance-weights would be manufactured for the use of buyers, and others for the use of sellers; and thus there would naturally be a tendency to the crystallization of a heavier and a lighter norm.
210. As for my assumption that the departures of the single weights from their virtual standards conform to the probability curve, it was only adopted as a ready way of imparting definiteness to the problem. Rich as is the store of data given by Petrie, it is insufficient, apparently, for determining the true law of those departures. If the workmen were sufficiently skillful (as I believe they would be) the departures would follow the probability curve. But if they were unskillful, it would be desirable to ascertain by what process the weights were made. The weights, being of stone, are not loaded; so that the adjustment was made by grinding, exclusively. Did the workman, then, have a balance by his side, or did he finish the weight by guesswork? In the latter case, inspection (and some sort of inspection there must, in this case, have been) would reject all weights outside a certain »tolerance,« as it is called in coinage. Those that were too light would have to be thrown away. They would lie in a heap, until they reappeared to deceive a future archeologist. Petrie's weights, however, are somewhat heavier, not lighter, than independent evidence would lead us to believe the ket to have been. Those that were too heavy would be reground, but would for the most part still be rather heavier than the standard. The consequence would be that the [error] curve would be cut down vertically at two ordinates (equally distant, perhaps, from the standard), while the ordinate of its maximum would be at the right of that of the standard. If the workman had a balance at hand, and frequently used it during the process of adjustment, the form of the error-curve would depend upon the construction of the balance. If it were like a modern balance, so as to show, not only that one mass is greater than another, but also whether it is much or little greater, the workman would keep in one pan a weight of the maximum value that he proposed to himself as permissible for the weight he was making; and in all his successive grindings would be aiming at that. The consequence would be a curve  concave upwards and stopping abruptly at its maximum ordinate: a form easily manageable by a slight modification of the method of least squares. But most of the balances shown upon the Egyptian monuments are provided with stops or other contrivances which would be needless if the balances were not top-heavy. Such balances, working automatically, are in use in all the mints of the civilized world, for throwing out light and heavy coins. Now a top-heavy balance will not show that two weights are equal, otherwise than by remaining with either end down which may be down. It only shows when, a weight being already in one pan, a decidedly heavier weight is placed in the other. The workman using such a balance would have no warning that he was approaching the limit, and would be unable to aim at any definite value, but (being, as we are supposing, devoid of skill), would have to grind away blindly, trying his weight every time he had ground off about as much as the whole range of variation which he proposed to allow himself. If he always ground off precisely the same amounts between successive tryings of his weight, he would be just as likely to grind below his maximum by any one fraction of the amount taken off at a grinding as by any other; so that his error curve would be a horizontal line cut off by vertical ordinates; thus, . But since there would be a variability in the amount taken off between the trials, the curve would show a contrary flexure; thus,
    
It must be admitted that the distribution of Petrie's kets is suggestive of this sort of curve, or rather of a modification of it due to a middling degree of skill.
211. I hope this long digression (which will be referred to with some interest when we come to study the theory of errors) will not have caused the reader to forget that we were engaged in tracing out some of the consequences of understanding the term »natural,« or »real,« class to mean a class the existence of whose members is due to a common and peculiar final cause. It is, as I was saying, a widespread error to think that a »final cause« is necessarily a purpose. A purpose is merely that form of final cause which is most familiar to our experience. The signification of the phrase »final cause« must be determined by its use in the statement of Aristotle 1) that all causation divides into two grand branches, the efficient, or forceful; and the ideal, or final. If we are to conserve the truth of that statement, we must understand by final causation that mode of bringing facts about according to which a general description of result is made to come about, quite irrespective of any compulsion for it to come about in this or that particular way; although the means may be adapted to the end. The general result may be brought about at one time in one way, and at another time in another way. Final causation does not determine in what particular way it is to be brought about, but only that the result shall have a certain general character.